\endinput
%------------------------------------------------------------------------------%
\begin{question}
\end{question}

%------------------------------------------------------------------------------%
\begin{resolution}{Solução}
\end{resolution}

%------------------------------------------------------------------------------%
\endinput



\begin{frame}{Exemplo -- SPI}
Considere
\[
\begin{cases}
x+2y=3\\
2x+4y=6.
\end{cases}
\]
\medskip

As duas equações representam a mesma reta.
\medskip

A segunda equação é o dobro da primeira:
\[
2(x+2y)=2(3).
\]
\medskip

Portanto, existem infinitas soluções.
\[
x+2y=3.
\]
Tomando $y=t$:
\[
x=3-2t.
\]
\medskip

Assim,
\[
\boxed{
(x,y)=(3-2t,t),\quad t\in\mathbb{R}.
}
\]
\end{frame}

\begin{frame}{Exemplo -- SI}
Agora considere
\[
\begin{cases}
x+2y=3\\
2x+4y=7.
\end{cases}
\]
\medskip

A matriz dos coeficientes é a mesma:
\[
A=
\begin{pmatrix}
1&2\\
2&4
\end{pmatrix},
\]
portanto,
\[
\det(A)=0.
\]
\medskip

Entretanto, o sistema é incompatível.
\end{frame}

\begin{frame}{Exemplo -- SI}
Multiplicando a primeira equação por $2$:
\[
2x+4y=6.
\]
\medskip

Mas a segunda equação afirma:
\[
2x+4y=7.
\]
\medskip

Subtraindo:
\[
0=1,
\]
uma contradição.
\medskip

Portanto,
\[
\boxed{\text{o sistema não possui solução}}.
\]
\end{frame}

